The following result generalizes Gårding\textquotesingle s inequality.

\begin{Theorem}
\label{thm:agmon-boundary-constrained-coercivity}
Let $L$, $B_j$ and $a$ be as above. Assume that Agmon\textquotesingle s condition holds at each point of $\partial\Omega$. Then there exist constants $c_1$ and $c_2$ such that \eqref{eq:agmon-boundary-constrained-coercive-estimate} holds.
\end{Theorem}

We now address the question how \eqref{eq:agmon-shifted-constrained-variational-equation} is to be interpreted as an elliptic boundary-value problem. For this, we first need a regularity statement.

\begin{Theorem}
\label{thm:agmon-weak-variational-solution-regularity}
Assume that $\Omega$ and the coefficients of $L$ and the $B_j$ are sufficiently smooth. Assume in addition that $f\in L^2(\Omega)$. Then the solution $u$ of \eqref{eq:agmon-shifted-constrained-variational-equation} lies in $H^{2m}(\Omega)$.
\end{Theorem}

Next, we need a Green	extquotesingle s formula.

\begin{Theorem}
\label{thm:agmon-constrained-form-green-boundary-identity}
Let $L$ and $a$ be as above. Let $B_i(\bm{x},D)$, $i=1,\ldots,m$, be a Dirichlet system of order $m$. Assume that $\Omega$ and the coefficients of the operators involved are sufficiently smooth. Then there exist normal boundary-value operators $C_i$, of order $2m-1-\operatorname{ord}B_i$, such that, for all $u,v\in H^{2m}(\Omega)$, we have
\begin{equation}\label{eq:agmon-constrained-form-green-boundary-identity}
a(u,v)=\int_\Omega(Lu)v\,\mathrm{d}\bm{x}-\sum_{i=1}^m\int_{\partial\Omega}(C_iu)(B_iv)\,\mathrm{d}S.
\end{equation}

\end{Theorem}

The proof is completely analogous to that of Theorem~\ref{thm:elliptic-green-formula-adjoint-boundary-operators}. For $u\in H^{2m}(\Omega)$ and $f\in L^2(\Omega)$, equation \eqref{eq:agmon-shifted-constrained-variational-equation} now assumes the form
\begin{equation}\label{eq:agmon-shifted-form-natural-boundary-reduction}
\int_\Omega(Lu+\lambda u)v\,\mathrm{d}\bm{x}-\sum_{j=p+1}^m\int_{\partial\Omega}(C_ju)(B_jv)\,\mathrm{d}S=\int_\Omega fv\,\mathrm{d}\bm{x}.
\end{equation}
This identifies \eqref{eq:agmon-shifted-constrained-variational-equation} as the weak form of the elliptic boundary-value problem
\begin{equation}\label{eq:agmon-essential-natural-boundary-value-problem}
Lu+\lambda u=f,\ B_ju=0,\ j=1,\ldots,p,\ C_ju=0,\ j=p+1,\ldots,m.
\end{equation}


The first set of boundary conditions is called essential; they are directly imposed on $u$ in the weak formulation of the problem. The second set of boundary conditions is called ``natural\textquotesingle\textquotesingle; they are not imposed explicitly, but arise from an integration by parts just like Neumann\textquotesingle s condition in Section~\ref{subsec:neumann-problem}.
